% Generated draft: requires mathematical and source review.
\documentclass[11pt]{article}
\usepackage[margin=25mm]{geometry}
\usepackage{fontspec}
\setmainfont{Latin Modern Roman}
\usepackage{amsmath,amssymb,mathtools,unicode-math}
\setmathfont{Latin Modern Math}
\usepackage{enumitem,needspace,xurl,hyperref}
\setlength{\parindent}{0pt}
\setlength{\parskip}{5pt}
\setlength{\emergencystretch}{3em}
\newcommand{\N}{\mathbb{N}}
\newcommand{\Z}{\mathbb{Z}}
\newcommand{\Q}{\mathbb{Q}}
\newcommand{\R}{\mathbb{R}}
\newcommand{\C}{\mathbb{C}}
\newcommand{\F}{\mathbb{F}}
\newcommand{\A}{\mathbb{A}}
\newcommand{\PP}{\mathbb{P}}
\newcommand{\E}{\mathbb{E}}
\newcommand{\eps}{\varepsilon}
\newcommand{\dd}{\,\mathrm{d}}
\begin{document}
\section*{UnsolvedMath ID: 30006490}
\textbf{Source problem code:} catalogue entry
\subsection{Definitions and mathematical statement}
Let \(G\) be a finite nilpotent group. Here nilpotent means that there is a finite central series \(1=G_0\leq G_1\leq\cdots\leq G_c=G\), with each \(G_i\) normal in \(G\) and \(G_{i+1}/G_i\subseteq Z(G/G_i)\). The integral group ring \(\Z[G]\) consists of sums \(\sum_{g\in G}a_g g\), with \(a_g\in\Z\), with multiplication induced by the group law. Define its augmentation and normalized unit group by
\[
\varepsilon:\Z[G]\longrightarrow\Z,\qquad \varepsilon\!\left(\sum_{g\in G}a_g g\right)=\sum_{g\in G}a_g,
\qquad
U_1(\Z[G])=\{u\in\Z[G]^\times:\varepsilon(u)=1\}.
\]
Identify \(G\) with its canonical subgroup of \(U_1(\Z[G])\).

A normal complement to \(G\) in \(U_1(\Z[G])\) is a subgroup \(N\trianglelefteq U_1(\Z[G])\) such that
\[
N\cap G=\{1\},\qquad NG=U_1(\Z[G]).
\]
Equivalently, \(U_1(\Z[G])=N\rtimes G\) as an internal semidirect product. Such a complement is torsion-free if, for every \(n\in N\) and every integer \(m\geq1\), the equality \(n^m=1\) implies \(n=1\).

\begin{enumerate}
\item For every finite nilpotent group \(G\), does \(G\) have a normal complement in \(U_1(\Z[G])\)?
\item For every finite nilpotent group \(G\) for which a normal complement exists, does there exist a torsion-free normal complement in \(U_1(\Z[G])\)? The torsion-free complement need not be the initially given complement.
\end{enumerate}

The supplied cleaned statement explicitly restricts these questions to finite nilpotent groups. The original excerpt instead says ``finite (nilpotent) group,'' leaving uncertain whether a broader question for arbitrary finite groups was also intended. The formulation above retains the explicit nilpotent restriction; it does not settle that scope discrepancy. The excerpt attributes the questions to Problem~29 in its reference~[2], whose bibliographic identity is not supplied. As motivation, the excerpt states that positive answers yield a positive solution to the integral group-ring isomorphism problem for \(G\); this is retained as a source claim, without supplying a proof or resolving which of the two positive answers the wording requires.
\subsection{Short English statement}
For each finite nilpotent group \(G\), ask whether its canonical copy inside the augmentation-one units of \(\Z[G]\) has a normal complement. Whenever such a complement exists, ask whether one can choose a torsion-free complement. The source's parenthetical use of ``nilpotent'' leaves the scope of its original wording uncertain.
\subsection{Sources}
\begin{itemize}
\item Supplied local UnsolvedMath registry record, permanent ID \(30006490\), title: Normal Complements in Integral Group-Ring Units; exact source problem code: \texttt{catalogue entry}. Collection: Oberwolfach Reports Open Problems. No individual proposer is identified; the supplied proposed year is 2025.
\item Supplied source citation: Mini-Workshop: The Yang--Baxter Equation and Representations of Braid Groups (2026), Oberwolfach Reports, pp.~2633--2672, DOI \texttt{10.4171/owr/2025/49}. \url{https://doi.org/10.4171/owr/2025/49}. The supplied excerpt carries page number 2653 and recalls Problem~29 in its reference~[2]; the identity of that reference is absent. The citation's parenthetical year 2026 differs from the report designation 49/2025 and the supplied proposed year 2025.
\item Supplied literature source: Oberwolfach Report 49/2025. \url{https://doi.org/10.4171/owr/2025/49}.
\item Supplied literature source: Counterexample to the integral group-ring isomorphism problem. \url{https://annals.math.princeton.edu/articles/10997}. The record supplies no authors or detailed explanation of this reference's bearing on the nilpotent-group questions.
\item Literature assessment dated \texttt{2026-08-22}: the supplied curation labels the problem open and reports that normal-complement existence for finite nilpotent groups remains unknown in general, with positive results for special families only. This is supplied dated curation derived from the Oberwolfach material, not independently verified here; the record itself requests independent verification of time-sensitive status and cited claims.
\item The registry labels the statement \texttt{corrected\_verified} and reports that the official report was checked, the question isolated, and notation repaired. These are supplied verification claims, not an independent verification performed for this draft.
\end{itemize}
\end{document}
